\documentclass{article}
\usepackage{amsmath}
\usepackage{geometry}
\geometry{a4paper, margin=1in}

\title{STM32 X-CUBE-AI Integration Steps}
\author{Generated by ChatGPT}
\date{}

\begin{document}

\maketitle

\section*{Step-by-Step Process for STM32 X-CUBE-AI Project Integration}

\begin{enumerate}
    \item \textbf{Create a New Project in STM32CubeMX}
    \begin{itemize}
        \item Open STM32CubeMX.
        \item Create a new project for your STM32 device (e.g., STM32F401RE).
        \item Configure the necessary peripherals for the project, such as GPIO, USART for debug printing, and any other peripherals you need for the application.
        \item Under the \textit{Middleware} section, enable \textbf{X-CUBE-AI}.
        \item Set up the clock configuration and ensure the clock speed is set to the maximum allowable for your device.
    \end{itemize}

    \item \textbf{Add X-CUBE-AI Component}
    \begin{itemize}
        \item In the \textit{Pinout \& Configuration} tab, go to \textbf{Software Packs} $\rightarrow$ \textbf{X-CUBE-AI}.
        \item Enable \textbf{AI Libraries}. This will include the necessary libraries for the AI framework.
    \end{itemize}

    \item \textbf{Import the Neural Network Model}
    \begin{itemize}
        \item Click on the \textbf{X-CUBE-AI} button on the toolbar.
        \item Click on \textbf{New Network} and choose the desired network name (e.g., \texttt{network}).
        \item Click \textbf{Import} and browse to your Keras \texttt{.h5} or \texttt{.keras} model file that you want to integrate.
        \item After importing the model, configure the parameters such as the number of layers, nodes, and activations as necessary.
        \item Click on \textbf{Analyze} to ensure the model is compatible and fits within the memory constraints of your device.
        \item If successful, you should see a green checkmark.
    \end{itemize}

    \item \textbf{Generate the Code}
    \begin{itemize}
        \item Once the model analysis is complete, click on \textbf{Generate Code}.
        \item Select the appropriate toolchain for code generation (e.g., STM32CubeIDE).
        \item After code generation, open the project in the STM32CubeIDE.
    \end{itemize}

    \item \textbf{Modify the Generated Code as Needed}
    \begin{itemize}
        \item Go to the \texttt{app\_x-cube-ai.c} file and locate the functions \texttt{MX\_X\_CUBE\_AI\_Init()} and \texttt{MX\_X\_CUBE\_AI\_Process()}.
        \item Add any necessary debug printing or modifications in the \textit{USER CODE} sections.
        \item Ensure the UART is configured correctly for serial output to monitor the model's output.
    \end{itemize}

    \item \textbf{Configure and Build the Project}
    \begin{itemize}
        \item Build the project in STM32CubeIDE.
        \item Ensure there are no build errors related to missing files or dependencies.
    \end{itemize}

    \item \textbf{Set Up Validation on Target}
    \begin{itemize}
        \item In the \textbf{X-CUBE-AI} GUI, set up \textbf{Validate on Target} by selecting the appropriate communication settings for your device.
        \item Select the USART that is connected to the ST-LINK’s Virtual COM port.
        \item Enable the checkbox for \textbf{Automatic Compilation and Download}.
        \item Click on \textbf{Validate on Target}.
    \end{itemize}

    \item \textbf{Generate the AI Validation Project}
    \begin{itemize}
        \item A new temporary project will be generated, compiled, and flashed to the device.
        \item This temporary project will run the validation on the device and send the results back to the host PC.
        \item Review the validation results to confirm if the model behaves as expected on the device.
    \end{itemize}

    \item \textbf{Integrate the Model in Your Application}
    \begin{itemize}
        \item Once validated, you can integrate the generated network code (found in \texttt{network.c} and \texttt{network\_data.c}) into your main application.
        \item Use the \texttt{ai\_network\_run()} API to run inference with your inputs and outputs.
    \end{itemize}
\end{enumerate}

\end{document}
